\section{Introducción: la convergencia de la IA y la biología}

Esta clase es la última del semestre de primavera de 2025 de MIT 6.S191 (Introduction to Deep Learning) y la imparte, como invitada especial, Ava Amini, investigadora de Microsoft Research (MSR). Amini tiene una formación que abarca tanto la biología de laboratorio húmedo como las ciencias computacionales: durante la licenciatura y el doctorado profundizó en experimentos celulares e investigación con modelos animales, y más adelante combinó su interés por la biología con métodos computacionales, dedicándose a impulsar la frontera de las ciencias de la vida mediante técnicas de IA.

\begin{knowledgebox}{La misión de Microsoft Research (MSR)}
La misión central de MSR es "avanzar la frontera de la ciencia y la tecnología en beneficio de la humanidad" (advance the frontiers of science and technology to benefit humanity). Su investigación se desarrolla desde tres niveles de perspectiva: (1) investigación de frontera, fundamental y de orientación académica; (2) despliegue responsable y éticamente adecuado de la tecnología; (3) un impacto positivo en la salud humana y en el descubrimiento científico en las ciencias naturales. Dentro de MSR existe un equipo dedicado, Biomedical Machine Learning (BioML), del cual Amini es integrante central.
\end{knowledgebox}

Amini señala que la biología es, en esencia, un sistema extremadamente complejo que opera a escala nanométrica. Incluso una visualización científica de la membrana de una sola célula, ampliada un millón de veces, revela una riqueza estructural sorprendente: una enorme cantidad de moléculas de proteína que desempeñan en ella todo tipo de funciones. Esta complejidad representa a la vez un enorme desafío para la IA y una oportunidad sin precedentes.

En el curso 6.S191, los estudiantes ya han aprendido los dos grandes paradigmas del \textbf{modelado predictivo} (predictive modeling) y el \textbf{modelado generativo} (generative modeling). Cuando este marco se aplica al ámbito biológico, la pregunta central sigue siendo la predicción y la generación, pero la modalidad de los datos cambia de manera fundamental:

\begin{itemize}
    \item Ya no se trata de texto en lenguaje natural, sino del \textbf{lenguaje de secuencias de las moléculas biológicas} individuales
    \item Ya no se trata de imágenes de rostros o de automóviles, sino de \textbf{representaciones celulares}, \textbf{imágenes de tejidos} y \textbf{datos de pacientes}
    \item El objetivo de predicción se convierte en: ¿cuál es la función de una molécula biológica? ¿Cómo responde una célula a un fármaco? ¿Cuál es el resultado clínico de un paciente?
    \item El objetivo de generación se convierte en: dada una función objetivo, ¿es posible diseñar una molécula biológica que la realice?
\end{itemize}

\begin{importantbox}{El paradigma bidireccional central de la IA para la biología}
\textbf{Predicción hacia adelante}: dada una molécula biológica $\rightarrow$ predecir su función (clasificación, regresión).\\
\textbf{Diseño inverso}: dada una función objetivo $\rightarrow$ generar una molécula biológica que la satisfaga (modelado generativo).\\
Estas dos direcciones se complementan y forman en conjunto el circuito completo mediante el cual la IA potencia a la biología. Lo fundamental es que tanto la predicción como el diseño de la IA deben, al final, conectarse con el mundo natural mediante \textbf{validación experimental}; esto es lo que distingue a la IA para la biología de los campos puramente computacionales.
\end{importantbox}

El contenido central de esta clase se centra en \textbf{la aplicación de la IA generativa al diseño de proteínas}, en particular en cómo trasladar el Diffusion Model (modelo de difusión), este potente marco de modelado generativo, del dominio de las imágenes al dominio de datos discretos de las secuencias de proteínas, y presenta el modelo \textbf{EvoDiff}, desarrollado por el equipo de Amini.

\subsection{Resumen de la sección}

La IA para la biología es una de las direcciones de aplicación de la IA con mayor potencial. Los dos pilares centrales de este campo son el modelado predictivo y el modelado generativo, y su particularidad radica en que necesariamente debe conectarse con el mundo natural mediante experimentación. Esta clase se centra en el uso de la IA generativa, en particular del Diffusion Model, para diseñar moléculas de proteína con una función objetivo.

